% Universal reusable tensor-network constructions for TNLean.
%
% The commands in this file compose tn_core.tex into standard local tensors,
% chains, closures, and finite patches.  A document loads the core before this
% file and then builds all complete diagrams from these common atoms.

\makeatletter

% ---------- Standard tensor-network sites ----------

% A canonical-form gauge X_{j,q} has two independent sector indices in
% addition to its horizontal bond indices.  The sector legs meet the glyph at
% distinct boundary ports: Block carries j and Copy carries q.  In particular,
% attaching both ports to parallel buses never identifies the two buses.
\newcommand{\TN@sectorgauge}[3]{%
    \TN@insertion[label=above:$#3$]{#1}{#2}%
    \TN@horizontalports{#1}%
    \TN@vport{#1Block}{#1}{240}%
    \TN@vport{#1Copy}{#1}{300}%
}

% The inverse gauge reflects the two lower legs while retaining their roles.
% This is the planar order in the canonical-form figures: Copy is the
% lower-left port and Block is the lower-right port of X^{-1}.
\newcommand{\TN@inverseSectorgauge}[3]{%
    \TN@insertion[label=above:$#3$]{#1}{#2}%
    \TN@horizontalports{#1}%
    \TN@vport{#1Block}{#1}{300}%
    \TN@vport{#1Copy}{#1}{240}%
}

\newcommand{\TN@factorwithlegs}[3]{%
    \TN@factor[minimum width=0.72cm, minimum height=0.52cm,
        inner sep=1pt]{#1}{#2}{#3}%
    \TN@mpoports{#1}%
    \TN@popenport{#1N}{0,0.26}%
    \TN@popenport{#1S}{0,-0.26}%
}

\newcommand{\TN@openpverticalports}[2][]{%
    \TN@popenport[#1]{#2N}{0,\TN@physicalleglen}%
    \TN@popenport[#1]{#2S}{0,-\TN@physicalleglen}%
}

% An operator or density matrix has paired system ports on each side.  The
% state glyph indicates that the internal tensor network has been suppressed.
\newcommand{\TN@operatorstate}[3]{%
    \TN@state{#1}{#2}{#3}%
    \TN@vwestport{#1WUp}{#1}{0.32}%
    \TN@vwestport{#1WDown}{#1}{0.68}%
    \TN@veastport{#1EUp}{#1}{0.32}%
    \TN@veastport{#1EDown}{#1}{0.68}%
}

% An MPS tensor has west and east virtual ports and one north physical port.
% The atom itself carries no open legs, so it may be composed without erasing
% pre-drawn index lines.
\newcommand{\TN@mpssite}[3]{%
    \TN@tensor{#1}{#2}%
    \TN@horizontalports{#1}%
    \TN@pport{#1N}{#1}{north}%
    \if\relax\detokenize{#3}\relax\else
        \node at ($(#1.south)+(0,-0.28)$) {$#3$};%
    \fi
}

% An MPO or MPDO tensor has west and east virtual ports and north and south
% physical ports.
\newcommand{\TN@mposite}[3]{%
    \TN@tensor{#1}{#2}%
    \TN@mpoports{#1}%
    \if\relax\detokenize{#3}\relax\else
        \node[anchor=west] at ($(#1.south east)+(0.12,-0.18)$) {$#3$};%
    \fi
}

% A square-lattice PEPS tensor has four cardinal virtual indices in addition
% to the north-east physical index.  Its five named boundary ports form the
% compositional interface; the base atom has no fixed open stubs.
\newcommand{\TN@pepssite}[3]{%
    \TN@tensor{#1}{#2}%
    \TN@cardinalports{#1}%
    \TN@pport{#1P}{#1}{45}%
    \if\relax\detokenize{#3}\relax\else
        \node[anchor=north east] at ($(#1.south west)+(-0.10,-0.10)$) {$#3$};%
    \fi
}

% A doubled local tensor with its physical index contracted between the two
% layers.  The prefix names the upper and lower tensor nodes.
\newcommand{\TN@doublelayer}[3]{%
    \TN@tensor{#1U}{#2}%
    \TN@tensor{#1D}{#3}%
    \TN@mpoports{#1U}%
    \TN@mpoports{#1D}%
    \TN@pconnectports{#1US}{#1DN}%
}

% Two displayed tensor-product factors.  Tensor product is written explicitly;
% it is never represented by an index line.
\newcommand{\TN@factorpair}[5]{%
    \TN@factor{#1L}{#2}{#4}%
    \TN@factor{#1R}{#3}{#5}%
    \node at ($(#1L.center)!0.5!(#1R.center)$) {$\otimes$};%
}

% ---------- Oriented fusion maps ----------

% Every trivalent map is an orientation of one constructor.  Its ports are
% named by mathematical role: Combined is the trunk, while FactorOne and
% FactorTwo are the ordered branches.  The orientation changes only the
% placement of these roles; it never changes the box or branch fractions.
\newcommand{\TN@trivalentport}[4]{%
    \csname TN@#1port\endcsname{#2}{#3}{#4}%
}

\newcommand{\TN@trivalentbetweenport}[5]{%
    \csname TN@#1#2port\endcsname{#3}{#4}{#5}%
}

\newcommand{\TN@trivalentmap}[6][]{%
    \def\TN@trikind{#6}%
    \def\TN@trivirtual{v}%
    \def\TN@triphysical{p}%
    \def\TN@trimorphism{m}%
    \ifx\TN@trikind\TN@trivirtual\else
    \ifx\TN@trikind\TN@triphysical\else
    \ifx\TN@trikind\TN@trimorphism\else
        \PackageError{tnlean-tikz}{Unknown trivalent-map port type `#6'}
            {Use `v' for virtual, `p' for physical, or `m' for morphism.}%
        \def\TN@trikind{v}%
    \fi\fi\fi
    \TN@map[#1]{#2}{#3}{#4}%
    \def\TN@triorientation{#5}%
    \def\TN@triright{right}%
    \def\TN@trileft{left}%
    \def\TN@tridown{down}%
    \def\TN@triup{up}%
    \ifx\TN@triorientation\TN@triright
        \TN@trivalentport{\TN@trikind}{#2Combined}{#2}{west}%
        \TN@trivalentbetweenport{\TN@trikind}{east}{#2FactorOne}{#2}
            {\TN@branchfirstfraction}%
        \TN@trivalentbetweenport{\TN@trikind}{east}{#2FactorTwo}{#2}
            {\TN@branchsecondfraction}%
    \else\ifx\TN@triorientation\TN@trileft
        \TN@trivalentport{\TN@trikind}{#2Combined}{#2}{east}%
        \TN@trivalentbetweenport{\TN@trikind}{west}{#2FactorOne}{#2}
            {\TN@branchfirstfraction}%
        \TN@trivalentbetweenport{\TN@trikind}{west}{#2FactorTwo}{#2}
            {\TN@branchsecondfraction}%
    \else\ifx\TN@triorientation\TN@tridown
        \TN@trivalentport{\TN@trikind}{#2Combined}{#2}{south}%
        \TN@trivalentport{\TN@trikind}{#2FactorOne}{#2}{north west}%
        \TN@trivalentport{\TN@trikind}{#2FactorTwo}{#2}{north east}%
    \else\ifx\TN@triorientation\TN@triup
        \TN@trivalentport{\TN@trikind}{#2Combined}{#2}{north}%
        \TN@trivalentport{\TN@trikind}{#2FactorOne}{#2}{south west}%
        \TN@trivalentport{\TN@trikind}{#2FactorTwo}{#2}{south east}%
    \else
        \PackageError{tnlean-tikz}{Unknown trivalent-map orientation `#5'}
            {Use `right', `left', `down', or `up'.}%
    \fi\fi\fi\fi
    \coordinate (#2LabelTop) at (#2.north);%
    \coordinate (#2LabelBottom) at (#2.south);%
}

\newcommand{\TN@splitmap}[4][]{%
    \TN@trivalentmap[#1]{#2}{#3}{#4}{right}{v}%
    \TN@aliasport{#2Mid}{#2Combined}%
    \TN@aliasport{#2Up}{#2FactorOne}%
    \TN@aliasport{#2Down}{#2FactorTwo}%
}

\newcommand{\TN@mergemap}[4][]{%
    \TN@trivalentmap[#1]{#2}{#3}{#4}{left}{v}%
    \TN@aliasport{#2Mid}{#2Combined}%
    \TN@aliasport{#2Up}{#2FactorOne}%
    \TN@aliasport{#2Down}{#2FactorTwo}%
}

\newcommand{\TN@fusionmap}[4][]{%
    \TN@trivalentmap[#1]{#2}{#3}{#4}{down}{v}%
    \TN@aliasport{#2Out}{#2Combined}%
    \TN@aliasport{#2Left}{#2FactorOne}%
    \TN@aliasport{#2Right}{#2FactorTwo}%
}

\newcommand{\TN@cofusionmap}[4][]{%
    \TN@trivalentmap[#1]{#2}{#3}{#4}{up}{v}%
    \TN@aliasport{#2In}{#2Combined}%
    \TN@aliasport{#2Left}{#2FactorOne}%
    \TN@aliasport{#2Right}{#2FactorTwo}%
}

% An action tensor V_ax^y is a trivalent virtual map with role names inherited
% from the MPO action: an MPO block a acts on an incoming state block x and
% produces an outgoing state block y.  Its reflected partner has the same
% interface names.
\newcommand{\TN@actionmap}[4][]{%
    \TN@splitmap[#1]{#2}{#3}{#4}%
    \TN@aliasport{#2StateOut}{#2Combined}%
    \TN@aliasport{#2MPO}{#2FactorOne}%
    \TN@aliasport{#2StateIn}{#2FactorTwo}%
}

\newcommand{\TN@coactionmap}[4][]{%
    \TN@mergemap[#1]{#2}{#3}{#4}%
    \TN@aliasport{#2StateOut}{#2Combined}%
    \TN@aliasport{#2MPO}{#2FactorOne}%
    \TN@aliasport{#2StateIn}{#2FactorTwo}%
}

% A physical splitter has one incoming physical index and two ordered outgoing
% physical indices.  The merger is its vertically reflected interface.  These
% atoms cover blocking and refinement isometries without assigning a virtual
% type to any physical Hilbert-space index.
\newcommand{\TN@physicalsplitmap}[4][]{%
    \TN@trivalentmap[#1]{#2}{#3}{#4}{up}{p}%
    \TN@aliasport{#2Trunk}{#2Combined}%
    \TN@aliasport{#2Left}{#2FactorOne}%
    \TN@aliasport{#2Right}{#2FactorTwo}%
}

\newcommand{\TN@physicalmergemap}[4][]{%
    \TN@trivalentmap[#1]{#2}{#3}{#4}{down}{p}%
    \TN@aliasport{#2Trunk}{#2Combined}%
    \TN@aliasport{#2Left}{#2FactorOne}%
    \TN@aliasport{#2Right}{#2FactorTwo}%
}

% ---------- One-dimensional words and closures ----------

% A one-dimensional tensor atom carries named boundary ports.  The centre port
% is used only to place labels; every index line starts or ends at a boundary
% port.
\newcommand{\TN@declareMPOWordTensor}[2]{%
    \TN@mposite{#1}{#2}{}%
    \coordinate (#1C) at (#1.center);%
}

\newcommand{\TN@openMPOword}[6]{%
    \TN@declareMPOWordTensor{#1L}{(0,0)}
    \TN@declareMPOWordTensor{#1M}{(\TN@chainsep,0)}
    \TN@declareMPOWordTensor{#1R}{(\TN@chainright,0)}
    \TN@openpverticalports{#1L}
    \TN@openpverticalports{#1M}
    \TN@openpverticalports{#1R}
    \TN@label{#1LN}{(0,\TN@physlabeloff)}{#2}
    \TN@label{#1LS}{(0,-\TN@physlabeloff)}{#3}
    \TN@label{#1RN}{(0,\TN@physlabeloff)}{#4}
    \TN@label{#1RS}{(0,-\TN@physlabeloff)}{#5}
    \TN@vopenport{#1LW}{-0.8,0}
    \TN@vconnectports{#1LE}{#1MW}
    \TN@vopenport{#1ME}{0.62,0}
    \TN@vopenport{#1RW}{-0.62,0}
    \TN@vopenport{#1RE}{0.8,0}
    \node at (\TN@chainellipsisx,0) {$\cdots$};
    \node at (\TN@chainellipsisx,-1.02) {$#6$};%
}

% Trace-closed MPO word with open ket and bra indices.  The two shown sites and
% the omitted middle lie on the common cell pitch; the trace height belongs to
% the selected layout profile.
\newcommand{\TN@closedMPOword}[2]{%
    \TN@declareMPOWordTensor{#1L}{(0,0)}
    \TN@mposite{#1M}{(\TN@cellpitch,0)}{#2}%
    \coordinate (#1MC) at (#1M.center);%
    \TN@declareMPOWordTensor{#1R}{({3*\TN@cellpitch},0)}
    \TN@openpverticalports{#1L}
    \TN@openpverticalports{#1M}
    \TN@openpverticalports{#1R}
    \TN@vconnectports{#1LE}{#1MW}
    \TN@vopenport{#1ME}{\TN@virtualleglen,0}
    \TN@vopenport{#1RW}{-\TN@virtualleglen,0}
    \TN@annotationterminal{#1Omitted}
        {({2*\TN@cellpitch},0)}%
    \TN@label{#1Omitted}{(0,0)}{\cdots}% chktex 11
    \TN@vtraceportsbelow{#1LW}{#1RE}
        {\TN@traceclearance}{\TN@periodictraceheight}%
}

% ---------- Sector buses ----------

% A sector bus is a virtual index line with named endpoints.  Taps are virtual
% contraction ports placed at relative positions along the bus, so arbitrary
% port-bearing atoms can be attached without reconstructing the line.
\newcommand{\TN@vbus}[5][]{%
    \TN@vterminal{#2L}{#3}%
    \TN@vterminal{#2R}{#4}%
    \TN@vconnectports[#1]{#2L}{#2R}%
    \if\relax\detokenize{#5}\relax\else
        \node[tn label, anchor=east] at ($(#2L)+(-0.10,0)$) {$#5$};%
    \fi
}

% A parallel bus inherits the horizontal extent of a previously declared bus.
% Its displacement is signed, so the same construction applies above or below.
\newcommand{\TN@vparallelbus}[5][]{%
    \TN@vbus[#1]{#2}{($(#3L)+(0,#4)$)}{($(#3R)+(0,#4)$)}{#5}%
}

\newcommand{\TN@vbustap}[3]{%
    \TN@vterminal{#2}{($(#1L)!#3!(#1R)$)}%
    \TN@anonymousjunction{(#2)}%
}

% Projecting a named point onto a bus aligns a tap with an attached atom
% without repeating its horizontal coordinate in the complete diagram.
\newcommand{\TN@vbusprojecttap}[3]{%
    \TN@assertporttype{#3}{virtual}%
    \TN@vterminal{#2}{($(#1L)!(#3)!(#1R)$)}%
    \TN@anonymousjunction{(#2)}%
}

% ---------- Repeated finite constructions ----------

% Short factor words occur in the blocked-channel constructions.  One indexed
% constructor fixes their pitch and interfaces; the arity-specific public
% names below are semantic aliases rather than independent geometries.
\newcommand{\TN@factorChain}[4]{%
    \coordinate (#1Centre) at #2;%
    \foreach \TN@index in {1,...,#3} {% chktex 11
        \TNVirtualFactorSite{#1Site\TN@index}
            {($(#1Centre)+({(\TN@index-1)*\TN@cellpitch},0)$)}{#4}%
    }%
    \pgfmathtruncatemacro{\TN@lastbond}{#3-1}%
    \foreach \TN@index [evaluate=\TN@index as \TN@next using
            int(\TN@index+1)] in {1,...,\TN@lastbond} {% chktex 11
        \TN@vconnectports{#1Site\TN@index E}{#1Site\TN@next W}%
    }%
    \TN@vopenport{#1Site1W}{-\TN@virtualleglen,0}%
    \TN@vopenport{#1Site#3E}{\TN@virtualleglen,0}%
}

\newcommand{\TN@declareThreeSiteChain}[4]{%
    \TN@declareMPOWordTensor{#1one}{(1.00,0)}
    \TN@declareMPOWordTensor{#1two}{(2.00,0)}
    \TN@declareMPOWordTensor{#1three}{(3.00,0)}
    \TN@label[anchor=north]{#1one}
        {(0,-\TN@traceclearance-0.12)}{#2}
    \TN@label[anchor=north]{#1two}
        {(0,-\TN@traceclearance-0.12)}{#3}
    \TN@label[anchor=north]{#1three}
        {(0,-\TN@traceclearance-0.12)}{#4}%
}

\newcommand{\TN@connectThreeSiteChain}[1]{%
    \TN@vconnectports{#1oneE}{#1twoW}
    \TN@vconnectports{#1twoE}{#1threeW}%
}

\newcommand{\TN@closeThreeSiteChain}[1]{%
    \TN@vtraceportsbelow{#1oneW}{#1threeE}
        {\TN@traceclearance}{\TN@traceclearance}%
}

\newcommand{\TN@openThreeSitePhysicalPorts}[1]{%
    \TN@popenport{#1oneN}{0,0.46}
    \TN@popenport{#1twoN}{0,0.46}
    \TN@popenport{#1threeN}{0,0.46}%
}

\newcommand{\TN@boxedThreeSiteChain}[4]{%
    \TN@declareThreeSiteChain{#1}{#2}{#3}{#4}
    \TN@connectThreeSiteChain{#1}
    \TN@closeThreeSiteChain{#1}
    \TN@openThreeSitePhysicalPorts{#1}%
}

\newcommand{\TN@boxedThreeSiteInsertion}[5]{%
    \TN@declareThreeSiteChain{#1}{#2}{#3}{#4}
    \TN@insertion[label=above:$#5$]{#1X}{(1.50,0)}
    \TN@vport{#1XW}{#1X}{west}
    \TN@vport{#1XE}{#1X}{east}
    \TN@vconnectports{#1oneE}{#1XW}
    \TN@vconnectports{#1XE}{#1twoW}
    \TN@vconnectports{#1twoE}{#1threeW}
    \TN@closeThreeSiteChain{#1}
    \TN@openThreeSitePhysicalPorts{#1}%
}

\newcommand{\TN@boxedThreeSitePhysicalOne}[5]{%
    \TN@declareThreeSiteChain{#1}{#2}{#3}{#4}
    \TN@insertion[label=left:$#5$]{#1O}{(1.00,0.72)}
    \TN@pport{#1ON}{#1O}{north}
    \TN@pport{#1OS}{#1O}{south}
    \TN@connectThreeSiteChain{#1}
    \TN@closeThreeSiteChain{#1}
    \TN@pconnectports{#1oneN}{#1OS}
    \TN@popenport{#1ON}{0,0.36}
    \TN@popenport{#1twoN}{0,0.46}
    \TN@popenport{#1threeN}{0,0.46}%
}

\newcommand{\TN@boxedThreeSitePhysicalTwo}[5]{%
    \TN@declareThreeSiteChain{#1}{#2}{#3}{#4}
    \TN@insertion[label=left:$#5$]{#1O}{(2.00,0.72)}
    \TN@pport{#1ON}{#1O}{north}
    \TN@pport{#1OS}{#1O}{south}
    \TN@connectThreeSiteChain{#1}
    \TN@closeThreeSiteChain{#1}
    \TN@popenport{#1oneN}{0,0.46}
    \TN@pconnectports{#1twoN}{#1OS}
    \TN@popenport{#1ON}{0,0.36}
    \TN@popenport{#1threeN}{0,0.46}%
}

% ---------- Public compositional calculus ----------

% These commands are the sole author-facing vocabulary.  Blueprint chapters,
% slides, galleries, and tests use them; the private \TN@... implementation is
% confined to tex/tn.  Each semantic atom exposes long mathematical role names.

\newenvironment{TNDiagram}[1][normal]{%
    \begingroup
    \TNLayoutProfile{#1}%
    \begin{tikzpicture}[tn picture,
        node distance=\TN@cellpitch cm and \TN@cellpitch cm]%
}{%
    \end{tikzpicture}%
    \endgroup
}

\newsavebox{\TN@equationrowbox}
\newenvironment{TNEquationRow}[1][compact]{%
    \par\smallskip
    \begin{lrbox}{\TN@equationrowbox}%
    \begingroup
    \TNLayoutProfile{#1}%
    \ignorespaces
}{%
    \unskip
    \endgroup
    \end{lrbox}%
    \noindent\hfill\usebox{\TN@equationrowbox}\hfill\mbox{}\par
    \smallskip
}

% Several related identities share one relation column.  Each cell remains an
% intrinsic TNTerm, while the tabular structure aligns the relations without
% introducing diagram-specific coordinates or padding.
\newenvironment{TNEquationRows}[1][compact]{%
    \par\smallskip
    \begingroup
    \TNLayoutProfile{#1}%
    \noindent\hfill
    \begin{tabular}{@{}r@{}c@{}l@{}}%
}{%
    \end{tabular}%
    \hfill\mbox{}\par
    \endgroup
    \smallskip
}

\newcommand{\TNTerm}[2][compact]{%
    \begingroup
    \TNLayoutProfile{#1}%
    \begin{tikzpicture}[tn picture,
        node distance=\TN@cellpitch cm and \TN@cellpitch cm]#2\end{tikzpicture}%
    \endgroup
}

\newcommand{\TNRelation}[1]{%
    \hspace{\TN@relationgap}\ensuremath{#1}\hspace{\TN@relationgap}%
}

% An omission is part of the topology of a word, not a freely placed label.
\newcommand{\TNOmission}[2]{%
    \node[tn label,#2] (#1) {$\cdots$};% chktex 11
    \TNEventInternal{motif|picture=\the\TN@pictureid|name=#1|kind=omission}%
}

% The omitted column of a double-layer word retains one mark in each layer.
\newcommand{\TNStackedOmission}[2]{%
    \TNOmission{#1Upper}{#2}%
    \node[tn label,below=\TN@stackedseparation of #1Upper]
        (#1Lower) {$\cdots$};% chktex 11
    \coordinate (#1Centre) at ($(#1Upper)!0.5!(#1Lower)$);%
    \TNEventInternal{motif|picture=\the\TN@pictureid|name=#1|kind=stacked-omission}%
}

% Public symbolic pitch for intrinsically composed terms.  Its numerical value
% belongs to the selected layout profile.
\newcommand{\TNCellPitch}{\TN@cellpitch}

\newcommand{\TNPortAlias}[2]{\TN@aliasport{#1}{#2}}
\newcommand{\TNVirtualPort}[3]{\TN@vport{#1}{#2}{#3}}
\newcommand{\TNPhysicalPort}[3]{\TN@pport{#1}{#2}{#3}}
\newcommand{\TNMorphismPort}[3]{\TN@mport{#1}{#2}{#3}}
\newcommand{\TNVirtualTerminal}[2]{\TN@vterminal{#1}{#2}}
\newcommand{\TNPhysicalTerminal}[2]{\TN@pterminal{#1}{#2}}
\newcommand{\TNMorphismTerminal}[2]{\TN@mterminal{#1}{#2}}
\newcommand{\TNVirtualWestPort}[3]{\TN@vwestport{#1}{#2}{#3}}
\newcommand{\TNVirtualEastPort}[3]{\TN@veastport{#1}{#2}{#3}}
\newcommand{\TNVirtualNorthPort}[3]{\TN@vnorthport{#1}{#2}{#3}}
\newcommand{\TNVirtualSouthPort}[3]{\TN@vsouthport{#1}{#2}{#3}}
\newcommand{\TNPhysicalWestPort}[3]{\TN@pwestport{#1}{#2}{#3}}
\newcommand{\TNPhysicalEastPort}[3]{\TN@peastport{#1}{#2}{#3}}
\newcommand{\TNPhysicalNorthPort}[3]{\TN@pnorthport{#1}{#2}{#3}}
\newcommand{\TNPhysicalSouthPort}[3]{\TN@psouthport{#1}{#2}{#3}}
\newcommand{\TNConnectVirtual}[2]{\TN@vconnectports{#1}{#2}}
\newcommand{\TNConnectPhysical}[2]{\TN@pconnectports{#1}{#2}}
\newcommand{\TNConnectMorphism}[2]{\TN@mconnectports{#1}{#2}}
\newcommand{\TNConnectVirtualHV}[2]{\TN@vconnectportshv{#1}{#2}}
\newcommand{\TNConnectVirtualVH}[2]{\TN@vconnectportsvh{#1}{#2}}
\newcommand{\TNConnectPhysicalHV}[2]{\TN@pconnectportshv{#1}{#2}}
\newcommand{\TNConnectPhysicalVH}[2]{\TN@pconnectportsvh{#1}{#2}}
\newcommand{\TNConnectMorphismHV}[2]{\TN@mconnectportshv{#1}{#2}}
\newcommand{\TNConnectMorphismVH}[2]{\TN@mconnectportsvh{#1}{#2}}
\newcommand{\TNTraceVirtualBelow}[2]{%
    \TN@vtraceportsbelow{#1}{#2}{\TN@traceclearance}{\TN@traceclearance}%
}
\newcommand{\TNTraceVirtualAbove}[2]{%
    \TN@vtraceportsabove{#1}{#2}{\TN@traceclearance}{\TN@traceclearance}%
}
\newcommand{\TNTraceVirtualRight}[2]{%
    \TN@vtraceportsright{#1}{#2}{\TN@traceclearance}{\TN@traceclearance}%
}
\newcommand{\TNTracePhysicalBelow}[2]{%
    \TN@ptraceportsbelow{#1}{#2}{\TN@traceclearance}{\TN@traceclearance}%
}
\newcommand{\TNTracePhysicalAbove}[2]{%
    \TN@ptraceportsabove{#1}{#2}{\TN@traceclearance}{\TN@traceclearance}%
}
\newcommand{\TNTracePhysicalRight}[2]{%
    \TN@ptraceportsright{#1}{#2}{\TN@traceclearance}{\TN@traceclearance}%
}

\newcommand{\TNTensor}[2]{\TN@tensor{#1}{#2}}
\newcommand{\TNComponent}[3]{\TN@component{#1}{#2}{#3}}
\newcommand{\TNFactor}[3]{\TN@factor{#1}{#2}{#3}}
\newcommand{\TNMap}[3]{\TN@map{#1}{#2}{#3}}
\newcommand{\TNState}[3]{\TN@state{#1}{#2}{#3}}
\newcommand{\TNExpression}[3]{\TN@expression{#1}{#2}{#3}}
\newcommand{\TNInsertion}[3]{\TN@labeledinsertion{#1}{#2}{#3}}
\newcommand{\TNJunction}[2]{\TN@junction{#1}{#2}}
\newcommand{\TNOperatorState}[3]{%
    \TN@operatorstate{#1}{#2}{#3}%
    \TN@aliasport{#1WestKet}{#1WUp}%
    \TN@aliasport{#1WestBra}{#1WDown}%
    \TN@aliasport{#1EastKet}{#1EUp}%
    \TN@aliasport{#1EastBra}{#1EDown}%
}
\newcommand{\TNSectorGauge}[3]{%
    \TN@sectorgauge{#1}{#2}{#3}%
    \TN@aliasport{#1West}{#1W}%
    \TN@aliasport{#1East}{#1E}%
}
\newcommand{\TNInverseSectorGauge}[3]{%
    \TN@inverseSectorgauge{#1}{#2}{#3}%
    \TN@aliasport{#1West}{#1W}%
    \TN@aliasport{#1East}{#1E}%
}

% Labels attach to a named object or port.  The four directional forms use a
% profile-controlled clearance; the general form is reserved for irregular
% regions and takes a relative displacement rather than a page coordinate.
\newcommand{\TNLabelAbove}[2]{%
    \TN@label[anchor=south]{#1}{(0,\TN@labelclearance)}{#2}%
}
\newcommand{\TNLabelBelow}[2]{%
    \TN@label[anchor=north]{#1}{(0,-\TN@labelclearance)}{#2}%
}
\newcommand{\TNLabelLeft}[2]{%
    \TN@label[anchor=east]{#1}{(-\TN@labelclearance,0)}{#2}%
}
\newcommand{\TNLabelRight}[2]{%
    \TN@label[anchor=west]{#1}{(\TN@labelclearance,0)}{#2}%
}
\newcommand{\TNSelectedLabelAbove}[2]{%
    \TN@label[anchor=south,tn region label selected]
        {#1}{(0,\TN@labelclearance)}{#2}%
}
\newcommand{\TNSelectedLabelBelow}[2]{%
    \TN@label[anchor=north,tn region label selected]
        {#1}{(0,-\TN@labelclearance)}{#2}%
}
\newcommand{\TNSecondaryLabelAbove}[2]{%
    \TN@label[anchor=south,tn region label secondary]
        {#1}{(0,\TN@labelclearance)}{#2}%
}
\newcommand{\TNSecondaryLabelBelow}[2]{%
    \TN@label[anchor=north,tn region label secondary]
        {#1}{(0,-\TN@labelclearance)}{#2}%
}

% A panel gives a named origin to one stage of a comparison without exposing
% TikZ scopes to theorem diagrams.  Its origin is global placement; all local
% geometry remains the responsibility of the atoms and motifs it contains.
\newenvironment{TNPanel}[2]{%
    \coordinate (#1) at #2;%
    \begin{scope}[shift={#2}]%
}{%
    \end{scope}%
}

% Fitted semantic boundaries are shared presentation operations.  The caller
% supplies only the already named objects which constitute the region.
\newcommand{\TNGroupFit}[2]{%
    \node[tn grouping region, fit=#2] (#1) {};%
}
\newcommand{\TNFactorFit}[2]{%
    \node[tn factor region, fit=#2] (#1) {};%
}

% Regions and explanatory marks do not represent index contractions.
\newcommand{\TNGroupingRegion}[3]{%
    \TN@groupregion{#1}{#2}{#3}%
}
\newcommand{\TNGroupingRegionAbove}[3]{%
    \TN@groupregionabove{#1}{#2}{#3}%
}
\newcommand{\TNGroupingRegionSelectedBelow}[3]{%
    \TN@groupregionselectedbelow{#1}{#2}{#3}%
}
\newcommand{\TNMPSSite}[3]{%
    \TN@mpssite{#1}{#2}{#3}%
    \TN@aliasport{#1West}{#1W}%
    \TN@aliasport{#1East}{#1E}%
    \TN@aliasport{#1Ket}{#1N}%
}

\newcommand{\TNMPOSite}[3]{%
    \TN@mposite{#1}{#2}{#3}%
    \TN@aliasport{#1West}{#1W}%
    \TN@aliasport{#1East}{#1E}%
    \TN@aliasport{#1Ket}{#1N}%
    \TN@aliasport{#1Bra}{#1S}%
}

% A displayed factor with the same four external roles as an MPO site.  It is
% used when the source decomposes a tensor into named local factors rather than
% introducing a new tensor site.
\newcommand{\TNVirtualFactorSite}[3]{%
    \TN@factorwithlegs{#1}{#2}{#3}%
    \TN@aliasport{#1West}{#1W}%
    \TN@aliasport{#1East}{#1E}%
    \TN@aliasport{#1Ket}{#1N}%
    \TN@aliasport{#1Bra}{#1S}%
}

\newcommand{\TNRotatedMPOSite}[3]{%
    \TN@tensor{#1}{#2}%
    \TN@vport{#1North}{#1}{north}%
    \TN@vport{#1South}{#1}{south}%
    \TN@pport{#1Ket}{#1}{west}%
    \TN@pport{#1Bra}{#1}{east}%
    \if\relax\detokenize{#3}\relax\else
        \TN@label{#1}{(-0.30,0.24)}{#3}%
    \fi
}

\newcommand{\TNPEPSSite}[3]{%
    \TN@pepssite{#1}{#2}{#3}%
    \TN@aliasport{#1West}{#1W}%
    \TN@aliasport{#1East}{#1E}%
    \TN@aliasport{#1North}{#1N}%
    \TN@aliasport{#1South}{#1S}%
    \TN@aliasport{#1Ket}{#1P}%
}

\newcommand{\TNDoubleLayer}[3]{%
    \TN@doublelayer{#1}{#2}{#3}%
    \TN@aliasport{#1UpperWest}{#1UW}%
    \TN@aliasport{#1UpperEast}{#1UE}%
    \TN@aliasport{#1LowerWest}{#1DW}%
    \TN@aliasport{#1LowerEast}{#1DE}%
}
\newcommand{\TNFactorPair}[5]{\TN@factorpair{#1}{#2}{#3}{#4}{#5}}

\newcommand{\TNOpenVirtualWest}[1]{\TN@vopenport{#1}{-\TN@virtualleglen,0}}
\newcommand{\TNOpenVirtualEast}[1]{\TN@vopenport{#1}{\TN@virtualleglen,0}}
\newcommand{\TNOpenVirtualNorth}[1]{\TN@vopenport{#1}{0,\TN@virtualleglen}}
\newcommand{\TNOpenVirtualSouth}[1]{\TN@vopenport{#1}{0,-\TN@virtualleglen}}
\newcommand{\TNOpenPhysicalNorth}[1]{\TN@popenport{#1}{0,\TN@physicalleglen}}
\newcommand{\TNOpenPhysicalSouth}[1]{\TN@popenport{#1}{0,-\TN@physicalleglen}}
\newcommand{\TNOpenPhysicalWest}[1]{\TN@popenport{#1}{-\TN@physicalleglen,0}}
\newcommand{\TNOpenPhysicalEast}[1]{\TN@popenport{#1}{\TN@physicalleglen,0}}

\newcommand{\TNTrivalentMap}[6]{%
    \TN@trivalentmap{#1}{#2}{#3}{#4}{#5}%
    \if\relax\detokenize{#6}\relax\else
        \TN@label{#1LabelTop}{(0,0.22)}{#6}%
    \fi
}
\newcommand{\TNTrivalentMapRight}[5]{%
    \TNTrivalentMap{#1}{#2}{#3}{right}{#4}{#5}}
\newcommand{\TNTrivalentMapLeft}[5]{%
    \TNTrivalentMap{#1}{#2}{#3}{left}{#4}{#5}}
\newcommand{\TNTrivalentMapDown}[5]{%
    \TNTrivalentMap{#1}{#2}{#3}{down}{#4}{#5}}
\newcommand{\TNTrivalentMapUp}[5]{%
    \TNTrivalentMap{#1}{#2}{#3}{up}{#4}{#5}}
\newcommand{\TNSplitMap}[3]{\TN@splitmap{#1}{#2}{#3}}
\newcommand{\TNMergeMap}[3]{\TN@mergemap{#1}{#2}{#3}}
\newcommand{\TNFusionMap}[3]{\TN@fusionmap{#1}{#2}{#3}}
\newcommand{\TNCofusionMap}[3]{\TN@cofusionmap{#1}{#2}{#3}}
\newcommand{\TNActionMap}[3]{\TN@actionmap{#1}{#2}{#3}}
\newcommand{\TNCoactionMap}[3]{\TN@coactionmap{#1}{#2}{#3}}
\newcommand{\TNPhysicalSplitMap}[3]{\TN@physicalsplitmap{#1}{#2}{#3}}
\newcommand{\TNPhysicalMergeMap}[3]{\TN@physicalmergemap{#1}{#2}{#3}}

% A stacked MPO product contracts exactly UpperBra with LowerKet.
\newcommand{\TNStackedMPOProduct}[5]{%
    \TNPlacementDispatch{#2}{%
        \TNMPOSite{#1Upper}{#2}{#3}%
        \TNMPOSite{#1Lower}
            {below=\TN@stackedseparation of #1Upper}{#4}%
        \coordinate (#1Centre) at
            ($(#1Upper.center)!0.5!(#1Lower.center)$);%
    }{%
        \coordinate (#1Centre) at #2;%
        \TNMPOSite{#1Upper}{($(#1Centre)+(0,\TN@layerpitch)$)}{#3}%
        \TNMPOSite{#1Lower}{($(#1Centre)+(0,-\TN@layerpitch)$)}{#4}%
    }%
    \TNConnectPhysical{#1UpperBra}{#1LowerKet}%
    \if\relax\detokenize{#5}\relax\else
        \TN@label{#1Centre}{(0.36,0)}{#5}%
    \fi
}

% A weighted block uses the same stacked product and adds one virtual weight
% on the outgoing upper bond.  The weight is part of the motif's interface.
\newcommand{\TNWeightedMPOBlock}[6]{%
    \TNStackedMPOProduct{#1}{#2}{#3}{#4}{#5}%
    \TN@insertion[label=above:$#6$]
        {#1Weight}{($(#1UpperEast)+(\TN@virtualleglen,0)$)}%
    \TN@horizontalports{#1Weight}%
    \TNConnectVirtual{#1UpperEast}{#1WeightW}%
    \TN@aliasport{#1WeightedEast}{#1WeightE}%
}

% Purification has two distinct ancillary ports joined by a straight physical
% contraction.  The external upper leg is Ket and the lower leg is Bra.
\newcommand{\TNPurificationSite}[4]{%
    \TNPlaceObjectInternal{#1Centre}{#2}%
    \coordinate (#1Centre) at (#1CentrePlacement);%
    \TN@tensor{#1KetTensor}{($(#1Centre)+(0,\TN@layerpitch)$)}%
    \TN@tensor{#1BraTensor}{($(#1Centre)+(0,-\TN@layerpitch)$)}%
    \TN@horizontalports{#1KetTensor}%
    \TN@horizontalports{#1BraTensor}%
    \TN@pport{#1Ket}{#1KetTensor}{north}%
    \TN@pport{#1KetAncilla}{#1KetTensor}{south east}%
    \TN@pport{#1BraAncilla}{#1BraTensor}{north east}%
    \TN@pport{#1Bra}{#1BraTensor}{south}%
    \TNConnectPhysical{#1KetAncilla}{#1BraAncilla}%
    \TN@annotationterminal{#1AncillaMid}
        {($(#1KetAncilla)!0.5!(#1BraAncilla)$)}%
    \TN@label{#1KetTensor}{(-0.30,0.24)}{#3}%
    \TN@label{#1BraTensor}{(-0.30,-0.24)}{#4}%
    \TN@aliasport{#1KetWest}{#1KetTensorW}%
    \TN@aliasport{#1KetEast}{#1KetTensorE}%
    \TN@aliasport{#1BraWest}{#1BraTensorW}%
    \TN@aliasport{#1BraEast}{#1BraTensorE}%
}

% A compact trace cell is the common definition of M_alpha(X) and A_alpha(X).
\newcommand{\TNCompactTraceCell}[5]{%
    \TNPlaceObjectInternal{#1Centre}{#2}%
    \coordinate (#1Centre) at (#1CentrePlacement);%
    \TN@insertion[label=above:$#4$]
        {#1Insertion}{($(#1Centre)+(-\TN@cellpitch/2,0)$)}%
    \TN@horizontalports{#1Insertion}%
    \TNMPOSite{#1Site}{($(#1Centre)+(\TN@cellpitch/2,0)$)}{}%
    \TNConnectVirtual{#1InsertionE}{#1SiteWest}%
    \TNTraceVirtualBelow{#1InsertionW}{#1SiteEast}%
    \TN@label{#1Centre}{(0,-\TN@traceclearance-0.22)}{\tr}%
    \TN@label{#1Site}{(0.32,0.26)}{#3}%
    \if\relax\detokenize{#5}\relax\else
        \TN@label{#1Centre}
            {(0,\TN@layerpitch+\TN@labelclearance)}{#5}%
    \fi
}

% Horizontal and rotated vertical words use the common pitches and closures.
\newcommand{\TNHorizontalWord}[4]{\TN@openMPOword{#1}{#2}{#3}{#2}{#3}{#4}}
% A periodic MPO word has the same indexed endpoints as an open word, but its
% outer virtual ports are joined by the canonical rounded trace closure.
\newcommand{\TNIndexedPeriodicMPOWord}[7]{%
    \TN@closedMPOword{#1}{#7}%
    \TN@label{#1LN}{(0,\TN@physlabeloff)}{#2}%
    \TN@label{#1LS}{(0,-\TN@physlabeloff)}{#3}%
    \TN@label{#1RN}{(0,\TN@physlabeloff)}{#4}%
    \TN@label{#1RS}{(0,-\TN@physlabeloff)}{#5}%
    \TN@annotationterminal{#1Length}
        {($(#1Omitted)+(0,-\TN@traceclearance-\TN@labelclearance)$)}%
    \TN@label{#1Length}{(0,0)}{#6}%
}
\newcommand{\TNVerticalWord}[4]{%
    \coordinate (#1Centre) at #2;%
    \TNRotatedMPOSite{#1Top}{($(#1Centre)+(0,\TN@cellpitch)$)}{#3}%
    \TNRotatedMPOSite{#1Bottom}{($(#1Centre)+(0,-\TN@cellpitch)$)}{#3}%
    \TN@vconnectports{#1TopSouth}{#1BottomNorth}%
    \node at (#1Centre) {$\vdots$};%
    \TNOpenPhysicalWest{#1TopKet}%
    \TNOpenPhysicalEast{#1TopBra}%
    \TNOpenPhysicalWest{#1BottomKet}%
    \TNOpenPhysicalEast{#1BottomBra}%
    \if\relax\detokenize{#4}\relax\else
        \TNTraceVirtualRight{#1TopNorth}{#1BottomSouth}%
    \fi
}

\newcommand{\TNSectorBus}[4]{\TN@vbus{#1}{#2}{#3}{#4}}
\newcommand{\TNParallelSectorBus}[4]{\TN@vparallelbus{#1}{#2}{#3}{#4}}
\newcommand{\TNSectorBusTap}[3]{\TN@vbustap{#1}{#2}{#3}}
\newcommand{\TNSectorBusProjectTap}[3]{%
    \TN@vbusprojecttap{#1}{#2}{#3}%
}
\newcommand{\TNFactorChainTwo}[3]{\TN@factorChain{#1}{#2}{2}{#3}}
\newcommand{\TNFactorChainThree}[3]{\TN@factorChain{#1}{#2}{3}{#3}}
\newcommand{\TNFactorChainFour}[3]{\TN@factorChain{#1}{#2}{4}{#3}}

\newcommand{\TNThreeSiteChain}[4]{\TN@boxedThreeSiteChain{#1}{#2}{#3}{#4}}
\newcommand{\TNThreeSiteInsertion}[5]{%
    \TN@boxedThreeSiteInsertion{#1}{#2}{#3}{#4}{#5}%
}
\newcommand{\TNThreeSitePhysicalFirst}[5]{%
    \TN@boxedThreeSitePhysicalOne{#1}{#2}{#3}{#4}{#5}%
}
\newcommand{\TNThreeSitePhysicalSecond}[5]{%
    \TN@boxedThreeSitePhysicalTwo{#1}{#2}{#3}{#4}{#5}%
}
% A four-bond gauge star has one fixed geometry in every client.  The gauge
% labels face away from the central tensor, leaving every incidence visible.
\newcommand{\TNFourBondGaugeStar}[6]{%
    \coordinate (#1Centre) at #2;%
    \TNTensor{#1Tensor}{(#1Centre)}%
    \TNPhysicalPort{#1Physical}{#1Tensor}{north}%
    \TNPhysicalTerminal{#1PhysicalOpen}
        {($(#1Physical)+(0,\TN@physicalleglen)$)}%
    \TNConnectPhysical{#1Physical}{#1PhysicalOpen}%
    \TNInsertion{#1WestInsertion}
        {($(#1Centre)+(-\TN@gaugeradius,0)$)}{}%
    \TNInsertion{#1EastInsertion}
        {($(#1Centre)+(\TN@gaugeradius,0)$)}{}%
    \TNInsertion{#1SouthWestInsertion}
        {($(#1Centre)+(-\TN@gaugediagonal,-\TN@gaugeradius)$)}{}%
    \TNInsertion{#1SouthEastInsertion}
        {($(#1Centre)+(\TN@gaugediagonal,-\TN@gaugeradius)$)}{}%
    \TNVirtualPort{#1TensorWest}{#1Tensor}{west}%
    \TNVirtualPort{#1TensorEast}{#1Tensor}{east}%
    \TNVirtualPort{#1TensorSouthWest}{#1Tensor}{225}%
    \TNVirtualPort{#1TensorSouthEast}{#1Tensor}{315}%
    \TNVirtualPort{#1WestInner}{#1WestInsertion}{east}%
    \TNVirtualPort{#1WestOuter}{#1WestInsertion}{west}%
    \TNVirtualPort{#1EastInner}{#1EastInsertion}{west}%
    \TNVirtualPort{#1EastOuter}{#1EastInsertion}{east}%
    \TNVirtualPort{#1SouthWestInner}{#1SouthWestInsertion}{45}%
    \TNVirtualPort{#1SouthWestOuter}{#1SouthWestInsertion}{225}%
    \TNVirtualPort{#1SouthEastInner}{#1SouthEastInsertion}{135}%
    \TNVirtualPort{#1SouthEastOuter}{#1SouthEastInsertion}{315}%
    \TNConnectVirtual{#1WestInner}{#1TensorWest}%
    \TNConnectVirtual{#1TensorEast}{#1EastInner}%
    \TNConnectVirtual{#1TensorSouthWest}{#1SouthWestInner}%
    \TNConnectVirtual{#1TensorSouthEast}{#1SouthEastInner}%
    \TNOpenVirtualWest{#1WestOuter}%
    \TNOpenVirtualEast{#1EastOuter}%
    \TNOpenVirtualSouth{#1SouthWestOuter}%
    \TNOpenVirtualSouth{#1SouthEastOuter}%
    \TN@annotationterminal{#1Caption}
        {($(#1Centre)+(0,-\TN@gaugeradius-\TN@virtualleglen)$)}%
    \TN@label[anchor=south]{#1WestInsertion.north}
        {(0,\TN@labelclearance)}{#3}%
    \TN@label[anchor=south]{#1EastInsertion.north}
        {(0,\TN@labelclearance)}{#4}%
    \TN@label[anchor=north east]{#1SouthWestInsertion.south west}
        {(-\TN@labelclearance,-\TN@labelclearance)}{#5}%
    \TN@label[anchor=north west]{#1SouthEastInsertion.south east}
        {(\TN@labelclearance,-\TN@labelclearance)}{#6}%
}

% Cardinal version used for a translation-invariant square PEPS tensor.
\newcommand{\TNCardinalGaugeCross}[6]{%
    \coordinate (#1Centre) at #2;%
    \TNTensor{#1Tensor}{(#1Centre)}%
    \TNPhysicalPort{#1Physical}{#1Tensor}{45}%
    \TNOpenPhysicalNorth{#1Physical}%
    \TNInsertion{#1WestInsertion}
        {($(#1Centre)+(-\TN@gaugeradius,0)$)}{}%
    \TNInsertion{#1EastInsertion}
        {($(#1Centre)+(\TN@gaugeradius,0)$)}{}%
    \TNInsertion{#1SouthInsertion}
        {($(#1Centre)+(0,-\TN@gaugeradius)$)}{}%
    \TNInsertion{#1NorthInsertion}
        {($(#1Centre)+(0,\TN@gaugeradius)$)}{}%
    \TNVirtualPort{#1TensorWest}{#1Tensor}{west}%
    \TNVirtualPort{#1TensorEast}{#1Tensor}{east}%
    \TNVirtualPort{#1TensorSouth}{#1Tensor}{south}%
    \TNVirtualPort{#1TensorNorth}{#1Tensor}{north}%
    \TNVirtualPort{#1WestInner}{#1WestInsertion}{east}%
    \TNVirtualPort{#1WestOuter}{#1WestInsertion}{west}%
    \TNVirtualPort{#1EastInner}{#1EastInsertion}{west}%
    \TNVirtualPort{#1EastOuter}{#1EastInsertion}{east}%
    \TNVirtualPort{#1SouthInner}{#1SouthInsertion}{north}%
    \TNVirtualPort{#1SouthOuter}{#1SouthInsertion}{south}%
    \TNVirtualPort{#1NorthInner}{#1NorthInsertion}{south}%
    \TNVirtualPort{#1NorthOuter}{#1NorthInsertion}{north}%
    \TNConnectVirtual{#1WestInner}{#1TensorWest}%
    \TNConnectVirtual{#1TensorEast}{#1EastInner}%
    \TNConnectVirtual{#1SouthInner}{#1TensorSouth}%
    \TNConnectVirtual{#1TensorNorth}{#1NorthInner}%
    \TNOpenVirtualWest{#1WestOuter}%
    \TNOpenVirtualEast{#1EastOuter}%
    \TNOpenVirtualSouth{#1SouthOuter}%
    \TNOpenVirtualNorth{#1NorthOuter}%
    \TN@annotationterminal{#1Caption}
        {($(#1Centre)+(0,-\TN@gaugeradius-\TN@virtualleglen)$)}%
    \TN@label[anchor=south]{#1WestInsertion.north}
        {(0,\TN@labelclearance)}{#3}%
    \TN@label[anchor=south]{#1EastInsertion.north}
        {(0,\TN@labelclearance)}{#4}%
    \TN@label[anchor=east]{#1SouthInsertion.west}
        {(-\TN@labelclearance,0)}{#5}%
    \TN@label[anchor=south]{#1NorthInsertion.north}
        {(0,\TN@labelclearance)}{#6}%
}

% A periodic MPS word owns its common pitch, omitted-site gap, trace closure,
% and trace-aware site-label baseline.
\newcommand{\TNCyclicMPSWord}[4]{%
    \coordinate (#1Centre) at (0,0);%
    \TNMPSSite{#1One}{(0,0)}{}%
    \TNMPSSite{#1Two}{(\TN@chainsep,0)}{}%
    \TNMPSSite{#1Penultimate}{(\TN@chainright-\TN@chainsep,0)}{}%
    \TNMPSSite{#1Last}{(\TN@chainright,0)}{}%
    \foreach \name/\index in
        {One/i_1,Two/i_2,Penultimate/i_{#3-1},Last/i_{#3}} {%
        \TNOpenPhysicalNorth{#1\name Ket}%
        \TNLabelAbove{#1\name Ket}{\index}%
        \TN@label[anchor=north]{#1\name.south}
            {(0,-\TN@labelclearance)}{#2}%
    }%
    \TNConnectVirtual{#1OneEast}{#1TwoWest}%
    \TNOpenVirtualEast{#1TwoEast}%
    \TNOpenVirtualWest{#1PenultimateWest}%
    \TNConnectVirtual{#1PenultimateEast}{#1LastWest}%
    \TNTraceVirtualBelow{#1OneWest}{#1LastEast}%
    \TN@annotationterminal{#1Omitted}{(\TN@chainellipsisx,0)}%
    \TN@annotationterminal{#1TraceLabel}
        {(\TN@chainellipsisx,-\TN@traceclearance-\TN@annotationclearance)}%
    \TN@label{#1Omitted}{(0,0)}{\cdots}% chktex 11
    \TN@label[anchor=north]{#1TraceLabel}{(0,0)}{#4}%
}

% Gauge conjugation of one MPS tensor is shared by blueprint and slides.
\newcommand{\TNGaugeConjugatedMPSSite}[5]{%
    \TNPlacementDispatch{#2}{%
        \TNMPSSite{#1Site}{#2}{#3}%
        \coordinate (#1Centre) at (#1Site.center);%
        \TNInsertion{#1Left}{left=of #1Site}{#4}%
        \TNInsertion{#1Right}{right=of #1Site}{#5}%
    }{%
        \coordinate (#1Centre) at #2;%
        \TNMPSSite{#1Site}{(#1Centre)}{#3}%
        \TNInsertion{#1Left}
            {($(#1Centre)+(-\TN@cellpitch,0)$)}{#4}%
        \TNInsertion{#1Right}
            {($(#1Centre)+(\TN@cellpitch,0)$)}{#5}%
    }%
    \TNVirtualPort{#1LeftWest}{#1Left}{west}%
    \TNVirtualPort{#1LeftEast}{#1Left}{east}%
    \TNVirtualPort{#1RightWest}{#1Right}{west}%
    \TNVirtualPort{#1RightEast}{#1Right}{east}%
    \TNOpenVirtualWest{#1LeftWest}%
    \TNConnectVirtual{#1LeftEast}{#1SiteWest}%
    \TNConnectVirtual{#1SiteEast}{#1RightWest}%
    \TNOpenVirtualEast{#1RightEast}%
    \TNOpenPhysicalNorth{#1SiteKet}%
}

\newcommand{\TNMorphismArrow}[4]{\TN@labeledmorphism{#1}{#2}{#3}{#4}}
\newcommand{\TNComparisonArrow}[4]{\TN@labeledcomparison{#1}{#2}{#3}{#4}}

% Two levels of the standard binary fusion tree.  All ports and branches are
% inherited from the unique trivalent-map constructor.
\newcommand{\TNBinaryFusionTree}[6]{%
    \coordinate (#1Centre) at #2;%
    \TNFusionMap{#1Lower}{($(#1Centre)+(0,-\TN@layerpitch)$)}{#6}%
    \TNFusionMap{#1Upper}{($(#1Centre)+(0,\TN@layerpitch)$)}{#6}%
    \TNConnectVirtual{#1UpperCombined}{#1LowerFactorOne}%
    \TN@label{#1UpperFactorOne}{(-0.22,0.20)}{#3}%
    \TN@label{#1UpperFactorTwo}{(0.22,0.20)}{#4}%
    \TN@label{#1LowerFactorTwo}{(0.22,0.20)}{#5}%
}

\makeatother
